\section{Related work}
\label{sec:related-work}

\paragraph{Separating systems and antichains.}
Separating systems on a finite ground set are classical
\citep{Katona1966,Dickson1969}.  In an ordinary separating system, the
incidence signatures of distinct elements need only be distinct.  Complete separation
requires, for every ordered pair \((e,f)\), a test containing \(e\) and
avoiding \(f\).  Equivalently, all incidence signatures must be pairwise
incomparable, connecting the problem directly with Sperner theory and minimal
completely separating systems \citep{Dickson1969,Spencer1970}.  The signature
argument in Proposition~\ref{prop:signature-antichain} is precisely this classical
equivalence after the admissible tests have been restricted to Hamiltonian
cycles.

\paragraph{Graph-structured separators.}
Graph-specific separating systems have a longer history; \citet{Cai1984}, for
example, studied separating systems of graphs.  We cite that work only as
related graph-specific context and do not assert that its definition coincides
with \(\hsep\).  A substantial line of work asks for edge separation by paths
\citep{FalgasRavryEtAl2014,BaloghEtAl2016,BonamyEtAl2023,Letzter2024}.
These problems use the same or a closely related directed separation
condition, but admit paths rather than Hamiltonian cycles as tests.  Recent
cycle-separation work allows arbitrary cycles or subdivisions of \(K_4\)
\citep{BotlerNaia2025}; this is a much larger admissible family than the
Hamiltonian cycles of a fixed graph.  Cycle-identification problems, such as
those studied for hypercubes and tori by \citet{HonkalaKarpovskyLitsyn2003},
generally seek distinct nonempty signatures.  Distinctness is weaker than the
pairwise incomparability required by \(\hsep\).

\paragraph{Hamiltonicity and Barnette graphs.}
Hamiltonian cycles in cubic 3-connected bipartite planar graphs were studied
computationally and structurally by \citet{HoltonManvelMcKay1985}.  Because the
complement of a Hamiltonian cycle in a cubic graph is a perfect matching,
Observation~\ref{obs:matching-formulation} also places the present parameter near matching
formulations of Barnette's conjecture \citep{GorskySteinerWiederrecht2023}.
Broader perspectives on that conjecture are given by \citet{AltEtAl2016}.
Our finite census uses the planar-graph generation methods of
\citet{BrinkmannMcKay2007} and Plantri 5.8 \citep{plantri58}.

The underlying notions of complete separation, antichain signatures, and edge
separation by unrestricted paths or cycles are therefore not new.  The
specific restriction here is that every admissible test must be a Hamiltonian
cycle of the given graph.  To the best of our knowledge, strong edge separation
under this restriction has not previously been investigated; this cautious
assessment is not an absolute novelty claim.
